\subsection{Moral Emotions: Empathy, Guilt, and Shame}
\label{sec:5.2.1}
Guilt and shame are self-conscious emotions: they fire when a person judges their own action, or their whole character, against a standard they have already accepted. Empathy, treated at length in section 2.2.1, supplies the input those judgments often run on: recognizing that an action caused harm is easier once the harm is felt and not merely inferred. What that earlier treatment does not cover, and what belongs here, is what guilt and shame specifically add to moral judgment: not just information about another's state, but a self-directed penalty for one's own conduct.

Guilt attaches to an act ("I did something wrong"); shame attaches to the self ("I am something wrong"). Both motivate correction, but they motivate it differently — guilt tends to produce reparative behavior, shame tends to produce withdrawal or defensiveness. An AI system built to learn from an ethical misstep could use a mechanism modeled on guilt specifically, since a generic negative-feedback signal will not do, if the goal is behavior change and not mere error correction: a feedback loop that penalizes the action taken, and passes no global judgment on the system's worth, has a chance of producing the corrective behavior guilt produces in people. Shame produces paralysis.

Building either into an AI system runs into the same problem twice over: a mechanism modeled on a human emotion needs an account of what that emotion is doing computationally, which is a harder question than what it feels like from the inside, and the ethical stakes of simulating an emotion capable of something like suffering are not smaller than the stakes of getting the underlying mechanism wrong. An autonomous vehicle that treats a near-miss with a pedestrian as guilt-relevant — weighting it in future route and speed decisions more heavily than a generic cost function would — is a plausible, bounded example of the first without committing to any claim about the second.

A separate strand worth naming here: guilt and shame in humans are entangled with mortality. Terror management theory holds that awareness of death shapes moral behavior by making people cling harder to worldviews and standards that offer a sense of continuity \autocite{greenberg1986causes}. An AI system has no comparable stake in its own continuation to ground an equivalent mechanism, which cuts both ways: it removes one lever that produces some prosocial behavior in humans, and it also removes a source of self-interested distortion. The more defensible design target is not to manufacture an artificial fear of death to replicate the human mechanism, but to build long-term, future-oriented weighting — the willingness to value consequences that fall well outside any single interaction — directly, without needing mortality salience to motivate it.
